% ---- score2 standard preamble (from: python3 score2_tools.py preamble) ----
\documentclass[10pt]{extarticle}
\usepackage[a4paper,margin=1.5cm]{geometry}
\usepackage[T1]{fontenc}
\usepackage{amsmath,amssymb}
\usepackage{graphicx}
\usepackage[export]{adjustbox}
\usepackage{multicol}
\usepackage{enumitem}
\usepackage{booktabs,array,tabularx}
\setlength{\columnsep}{0.7cm}
\setlength{\parindent}{0pt}
\setlength{\parskip}{3pt}
\pagestyle{empty}
\makeatletter
% figure registry (figdefs.tex is generated by the extract step)
\newcommand{\qfigdef}[4]{\@namedef{qf@file@#1}{#2}\@namedef{qf@w@#1}{#3}\@namedef{qf@wide@#1}{#4}}
\newcommand{\qf@check}[1]{\@ifundefined{qf@file@#1}{\PackageError{score2}{Unknown figure id #1}{}}{}}
\newcommand{\qf@img}[4][]{\includegraphics[width=#2\dimexpr\@nameuse{qf@w@#4}pt\relax,max width=#3,max height=0.45\textheight,#1]{\@nameuse{qf@file@#4}}}
% Full-width questions: finalize wraps every question that contains a wide figure in \qfullbegin...\qfullend,
% so the WHOLE question (stem, figures, options) is set across the page in reading order, then the two columns
% resume. The question's head - its number, stem and first figure (\qfullsplit marks the end of the head) -
% is kept on one page, so a number is never left at a page foot with its figure overleaf; the rest may break
% across pages normally, so no big gaps are left. Without a \qfullsplit, the whole question is kept together
% if it is under 45% of a page. Never write these yourself.
\newif\ifqf@full
\newif\ifqf@boxed
\newsavebox{\qf@box}
\newcommand{\qf@open}{\global\qf@fulltrue\par\global\qf@boxedtrue
  \setbox\qf@box\vbox\bgroup\hsize\textwidth\linewidth\textwidth}
\newcommand{\qf@emit}[1]{\par\egroup\global\qf@boxedfalse
  \ifdim\dimexpr\ht\qf@box+\dp\qf@box\relax<#1\textheight\noindent\usebox{\qf@box}\par\else\unvbox\qf@box\fi}
\newcommand{\qfullsplit}{\ifqf@boxed\qf@emit{0.85}\fi}
\newcommand{\qf@close}{\ifqf@boxed\qf@emit{0.45}\fi\global\qf@fullfalse}
\newcommand{\qfullbegin}{\end{multicols}\qf@open}
\newcommand{\qfullend}{\qf@close\begin{multicols}{2}}
\newcommand{\qfullstart}{\qf@open}
\newcommand{\qfulldone}{\qf@close}
% \qfig[scale]{id}: block figure, centred, directly where it is written (size set by the tool).
% Fallback if a wide figure is not inside a full-width question: break out of the columns in place - no floats.
\newcommand{\qfig}[2][1]{\qf@check{#2}%
  \ifqf@full\qf@blk{#1}{#2}\else
  \ifnum\@nameuse{qf@wide@#2}=1 \expandafter\@firstoftwo\else\expandafter\@secondoftwo\fi
  {\end{multicols}\begin{center}\qf@img{#1}{\textwidth}{#2}\end{center}\begin{multicols}{2}}%
  {\qf@blk{#1}{#2}}\fi}
\newcommand{\qf@blk}[2]{\par\smallskip{\centering\qf@img{#1}{\linewidth}{#2}\par}\smallskip}
% \qfigi[scale]{id}: inline figure (inside option lists / tables / minipages); never breaks columns.
\newcommand{\qfigi}[2][1]{\qf@check{#2}\qf@img[valign=c]{#1}{\linewidth}{#2}}
\makeatother
% structure macros - lint reads these to check completeness and figure placement
\newcommand{\papersec}[1]{\par\bigskip{\centering\Large\bfseries PAPER-#1\par}\medskip}
\newcommand{\partsec}[2]{\par\medskip\noindent{\large\bfseries PART-#1: #2}\par\smallskip}
\newcommand{\sect}[1]{\par\smallskip\noindent\textbf{#1}\par}
\newcommand{\qn}[1]{\par\smallskip\noindent\textbf{#1.}~}
\newcommand{\paraq}[2]{\par\smallskip\noindent\textbf{Paragraph for Questions #1 and #2}\par}
\newenvironment{qoptions}{\begin{itemize}[leftmargin=1.8em,itemsep=1pt,topsep=2pt,label={}]}{\end{itemize}}
\newcommand{\opt}[2]{\item[(#1)] #2}
% figure registry (generated by score2_tools.py finalize)
\qfigdef{p1_phy_s1_q01_1}{Score-II - Test 33 (01 Apr 2024)-img/p1_phy_s1_q01_1.png}{164.4}{0}
\qfigdef{p1_phy_s2_q02_1}{Score-II - Test 33 (01 Apr 2024)-img/p1_phy_s2_q02_1.png}{104.3}{0}
\qfigdef{p1_phy_s2_q03_1}{Score-II - Test 33 (01 Apr 2024)-img/p1_phy_s2_q03_1.png}{137.0}{0}
\qfigdef{p1_phy_s2_q05_1}{Score-II - Test 33 (01 Apr 2024)-img/p1_phy_s2_q05_1.png}{156.0}{0}
\qfigdef{p1_phy_s3_q02_1}{Score-II - Test 33 (01 Apr 2024)-img/p1_phy_s3_q02_1.png}{41.5}{0}
\qfigdef{p1_phy_s3_q03_1}{Score-II - Test 33 (01 Apr 2024)-img/p1_phy_s3_q03_1.png}{60.5}{0}
\qfigdef{p1_phy_s3_q03_2}{Score-II - Test 33 (01 Apr 2024)-img/p1_phy_s3_q03_2.png}{64.1}{0}
\qfigdef{p1_phy_s3_q03_3}{Score-II - Test 33 (01 Apr 2024)-img/p1_phy_s3_q03_3.png}{39.0}{0}
\qfigdef{p1_phy_s3_q03_4}{Score-II - Test 33 (01 Apr 2024)-img/p1_phy_s3_q03_4.png}{77.9}{0}
\qfigdef{p1_chem_s1_q03_1}{Score-II - Test 33 (01 Apr 2024)-img/p1_chem_s1_q03_1.png}{210.4}{0}
\qfigdef{p1_chem_s1_q03_2}{Score-II - Test 33 (01 Apr 2024)-img/p1_chem_s1_q03_2.png}{215.0}{0}
\qfigdef{p1_chem_s1_q04_1}{Score-II - Test 33 (01 Apr 2024)-img/p1_chem_s1_q04_1.png}{208.9}{0}
\qfigdef{p1_chem_s1_q05_1}{Score-II - Test 33 (01 Apr 2024)-img/p1_chem_s1_q05_1.png}{146.7}{0}
\qfigdef{p1_chem_s2_q02_1}{Score-II - Test 33 (01 Apr 2024)-img/p1_chem_s2_q02_1.png}{240.8}{0}
\qfigdef{p1_chem_s2_q03_1}{Score-II - Test 33 (01 Apr 2024)-img/p1_chem_s2_q03_1.png}{151.3}{0}
\qfigdef{p1_chem_s2_q05_1}{Score-II - Test 33 (01 Apr 2024)-img/p1_chem_s2_q05_1.png}{134.0}{0}
\qfigdef{p1_chem_s2_q05_2}{Score-II - Test 33 (01 Apr 2024)-img/p1_chem_s2_q05_2.png}{158.5}{0}
\qfigdef{p1_chem_s3_q03_B}{Score-II - Test 33 (01 Apr 2024)-img/p1_chem_s3_q03_B.png}{104.8}{0}
\qfigdef{p1_chem_s3_q03_C}{Score-II - Test 33 (01 Apr 2024)-img/p1_chem_s3_q03_C.png}{102.6}{0}
\qfigdef{p1_chem_s3_q03_D}{Score-II - Test 33 (01 Apr 2024)-img/p1_chem_s3_q03_D.png}{108.2}{0}
\qfigdef{p2_phy_s1_q02_1}{Score-II - Test 33 (01 Apr 2024)-img/p2_phy_s1_q02_1.png}{111.5}{0}
\qfigdef{p2_phy_s1_q05_1}{Score-II - Test 33 (01 Apr 2024)-img/p2_phy_s1_q05_1.png}{101.3}{0}
\qfigdef{p2_phy_s3_q01_1}{Score-II - Test 33 (01 Apr 2024)-img/p2_phy_s3_q01_1.png}{84.2}{0}
\qfigdef{p2_phy_s3_q04_1}{Score-II - Test 33 (01 Apr 2024)-img/p2_phy_s3_q04_1.png}{153.3}{0}
\qfigdef{p2_phy_s3_q06_1}{Score-II - Test 33 (01 Apr 2024)-img/p2_phy_s3_q06_1.png}{106.0}{0}
\qfigdef{p2_phy_s3_q07_1}{Score-II - Test 33 (01 Apr 2024)-img/p2_phy_s3_q07_1.png}{42.3}{0}
\qfigdef{p2_phy_s3_q08_1}{Score-II - Test 33 (01 Apr 2024)-img/p2_phy_s3_q08_1.png}{189.2}{0}
\qfigdef{p2_phy_s3_q10_1}{Score-II - Test 33 (01 Apr 2024)-img/p2_phy_s3_q10_1.png}{105.8}{0}
\qfigdef{p2_chem_s1_q03_A}{Score-II - Test 33 (01 Apr 2024)-img/p2_chem_s1_q03_A.png}{73.9}{0}
\qfigdef{p2_chem_s1_q03_B}{Score-II - Test 33 (01 Apr 2024)-img/p2_chem_s1_q03_B.png}{70.6}{0}
\qfigdef{p2_chem_s1_q03_C}{Score-II - Test 33 (01 Apr 2024)-img/p2_chem_s1_q03_C.png}{75.4}{0}
\qfigdef{p2_chem_s1_q03_D}{Score-II - Test 33 (01 Apr 2024)-img/p2_chem_s1_q03_D.png}{75.1}{0}
\qfigdef{p2_chem_s1_q04_A}{Score-II - Test 33 (01 Apr 2024)-img/p2_chem_s1_q04_A.png}{40.4}{0}
\qfigdef{p2_chem_s1_q04_B}{Score-II - Test 33 (01 Apr 2024)-img/p2_chem_s1_q04_B.png}{45.6}{0}
\qfigdef{p2_chem_s1_q04_C}{Score-II - Test 33 (01 Apr 2024)-img/p2_chem_s1_q04_C.png}{40.6}{0}
\qfigdef{p2_chem_s1_q04_D}{Score-II - Test 33 (01 Apr 2024)-img/p2_chem_s1_q04_D.png}{40.2}{0}
\qfigdef{p2_chem_s2_para9_10_1}{Score-II - Test 33 (01 Apr 2024)-img/p2_chem_s2_para9_10_1.png}{245.2}{0}
\qfigdef{p2_chem_s2_q09_C}{Score-II - Test 33 (01 Apr 2024)-img/p2_chem_s2_q09_C.png}{150.9}{0}
\qfigdef{p2_chem_s2_q09_D}{Score-II - Test 33 (01 Apr 2024)-img/p2_chem_s2_q09_D.png}{157.0}{0}
\qfigdef{p2_chem_s2_q10_A}{Score-II - Test 33 (01 Apr 2024)-img/p2_chem_s2_q10_A.png}{116.6}{0}
\qfigdef{p2_chem_s2_q10_B}{Score-II - Test 33 (01 Apr 2024)-img/p2_chem_s2_q10_B.png}{115.3}{0}
\qfigdef{p2_chem_s2_q10_C}{Score-II - Test 33 (01 Apr 2024)-img/p2_chem_s2_q10_C.png}{117.8}{0}
\qfigdef{p2_chem_s2_q10_D}{Score-II - Test 33 (01 Apr 2024)-img/p2_chem_s2_q10_D.png}{120.5}{0}
\qfigdef{p2_chem_s2_q11_1}{Score-II - Test 33 (01 Apr 2024)-img/p2_chem_s2_q11_1.png}{122.2}{0}
\qfigdef{p2_chem_s3_q02_1}{Score-II - Test 33 (01 Apr 2024)-img/p2_chem_s3_q02_1.png}{176.6}{0}
\qfigdef{p2_chem_s3_q04_1}{Score-II - Test 33 (01 Apr 2024)-img/p2_chem_s3_q04_1.png}{188.9}{0}
\qfigdef{p2_chem_s3_q05_1}{Score-II - Test 33 (01 Apr 2024)-img/p2_chem_s3_q05_1.png}{142.8}{0}
\qfigdef{p2_chem_s3_q06_1}{Score-II - Test 33 (01 Apr 2024)-img/p2_chem_s3_q06_1.png}{174.4}{0}
\qfigdef{p2_chem_s3_q06_2}{Score-II - Test 33 (01 Apr 2024)-img/p2_chem_s3_q06_2.png}{231.9}{0}
\qfigdef{p2_chem_s3_q07_1}{Score-II - Test 33 (01 Apr 2024)-img/p2_chem_s3_q07_1.png}{30.3}{0}
\qfigdef{p2_math_s3_q05_1}{Score-II - Test 33 (01 Apr 2024)-img/p2_math_s3_q05_1.png}{148.4}{0}
\raggedcolumns
\makeatletter\@beginparpenalty=10000\makeatother
\begin{document}
\begin{center}{\Large\bfseries JEE (Main+Advanced) Score-II --- Test 33 (01 Apr 2024)}\end{center}
\begin{multicols}{2}

\papersec{1}
\partsec{1}{PHYSICS}
\sect{Section-I}

\qn{1} A very long uniformly rough horizontal conveyor belt is moved with constant velocity 6 m/s. A small coin is projected on belt from origin of a co-ordinate system as shown in figure. Which is fixed to the ground so that initially it starts moving with speed v = 8 m/s on the belt along y-axis. The coefficient of friction between coin and surface of belt is $\mu=0.6$, the width of belt d = 5 m, take g = 10 m/s$^2$.
\qfig{p1_phy_s1_q01_1}
\begin{qoptions}
\opt{A}{Coin moves with constant acceleration on the belt}
\opt{B}{Coin moves with non uniform acceleration on the belt}
\opt{C}{Coin leaves the belt with speed 5 m/s}
\opt{D}{Finally coin comes to rest on the belt}
\end{qoptions}

\qn{2} A satellite of mass m orbits around the Earth in an elliptical path of semi-major and semi-minor axis 2a and a respectively such that the Earth is located at one of the foci of the elliptical path. At an instant when the satellite is located at a distance of r , its velocity is u at an angle of $60^\circ$ with respect to the position vector of the satellite from the centre of the Earth. Consider only the gravitational force between the satellite and the Earth and assume that the Earth is a perfect sphere.\\
(Assume that mass of the satellite is negligible as compared to mass of the Earth)
\begin{qoptions}
\opt{A}{Orbital angular momentum of the satellite with respect to centre of the Earth is $\dfrac{\sqrt3}{2}mur$}
\opt{B}{Areal velocity of position vector of the satellite with respect to centre of the Earth is $\dfrac{\sqrt3}{4}ur$}
\opt{C}{Time period of orbital motion of the satellite around the Earth is $\dfrac{8\pi a^2}{\sqrt3\,ur}$}
\opt{D}{Time period of satellite's orbital motion is independent of mass of the Earth}
\end{qoptions}

\qn{3} Max Planck noted in 1899 the existence of a system of units based on the three fundamental constants G, c, and h. These constants are dimensionally independent in the sense that no combination is dimensionless and a length, a time, and a mass may be constructed from them.
\begin{qoptions}
\opt{A}{The length can be represented by formula $\sqrt{\dfrac{\hbar G}{c^3}}$}
\opt{B}{The length can be represented by formula $\sqrt{\dfrac{\hbar c}{G}}$}
\opt{C}{The dimensional formula for acceleration is $\hbar^{3/2}c^{5/2}G^{1/2}$}
\opt{D}{The dimensional formula for acceleration is $\hbar^{-1/2}c^{7/2}G^{-1/2}$}
\end{qoptions}

\qn{4} A stationary sound source emits a wave that is received by a stationary observer as $p=p_0\sin(kx-wt)$.
\begin{qoptions}
\opt{A}{If source moves towards the observer, he will perceive higher k as well as higher $\omega$}
\opt{B}{If source moves towards the observer, he will perceive lower k but higher $\omega$}
\opt{C}{If observer moves towards the source, he will perceive same k but higher $\omega$}
\opt{D}{If observer moves towards the source, he will perceive higher k as well as higher $\omega$}
\end{qoptions}

\qn{5} A parallel plate air capacitor is connected to a battery. If plates of the capacitor are pulled further apart, then which of the following statements are correct?
\begin{qoptions}
\opt{A}{Strength of electric field inside the capacitor remain unchanged, if battery is disconnected before pulling the plate}
\opt{B}{During the process, positive work is done by an external force applied to pull the plates whether battery is disconnected or it remain connected}
\opt{C}{Potential energy in the capacitor decreases if the battery remains connected during pulling the plates apart}
\opt{D}{Potential energy in the capacitor decreases if the battery is disconnected before pulling the plates apart}
\end{qoptions}

\sect{Section-III}

\qn{1} The main scale of a vernier calipers reads in millimeter and its vernier is divided into ten divisions, which coincide with seven divisions of main scale. Furthermore when a cylinder is tightly placed along its length between two jaws, it is observed that the zero of vernier scale lies just right to 33$^{\text{rd}}$ divisions of main scale and fourth division of vernier scale coincide with the one of the main scale division. Then the measured value is 0.83 K cm. Find the value of K.

\qn{2} A solid conducting sphere of radius `a' having charge q is surrounded by a concentric conducting spherical shell of inner radius 2a and outer radius 3a as shown in figure. The net charge on spherical shell is $-2q$. If the amount of heat produced when switch is closed is $x\left(\dfrac{q^2}{96\pi\varepsilon_0a}\right)$. Fill the value of x.
\qfig{p1_phy_s2_q02_1}

\qn{3} A long plank P is placed on a horizontal frictionless surface and connected to two springs of spring constant k and 2k respectively. A cylinder of mass m is placed on the plank. The plank is displaced towards right (positive direction) and released at t = 0 and it oscillates with amplitude A. The cylinder does not slip on the plank during the oscillation. The maximum friction force acting on the cylinder during the oscillation is $\dfrac{PkA}{7}$. Fill the value of P.
\qfig{p1_phy_s2_q03_1}

\qn{4} On a planet which has the same density as that of earth, they construct a frictionless inclined plane of inclination $53^\circ$ with the horizontal. When a block is released on the inclined plane, it is seen to accelerate at 1 m/s$^2$. What is the radius of planet (in $10^5$m ). Gravitational acceleration at surface of earth is 10 m/s$^2$, Radius of earth 6400 km.

\qn{5} Three balls are kept on a straight line on a smooth surface. The first ball is projected towards middle ball as shown. It collides with middle ball of mass m$_2$ which then collides with m$_3$ . We want m$_3$ to move with maximum possible velocity after one collision. What should be mass m$_2$ (in kg)? All collisions are head on and elastic.
\qfig{p1_phy_s2_q05_1}

\sect{Section-IV}

\qn{1} An ideal monatomic gas undergoes different types of processes which are described in Column $-$ I match the corresponding effects in column $-$ II. The letters have their usual meaning.\par\nobreak
\begin{center}\small\setlength{\tabcolsep}{3pt}
\begin{tabular}{|c|p{2.2cm}|c|p{3.2cm}|}
\hline
\multicolumn{2}{|c|}{Column$-$I} & \multicolumn{2}{c|}{Column$-$II} \\
\hline
(A) & $PV^{-2}$ = constant & (P) & If volume increases then temperature will also increase \\
\hline
(B) & $PV^{2}$ = constant & (Q) & For expansion, heat will have to be supplied to the gas \\
\hline
(C) & C = C$_v$ + 2R & (R) & If temperature increases then work done by gas is positive \\
\hline
(D) & C = C$_v$ $-$ 2R & (S) & If pressure increases the work done by gas is negative \\
\hline
\end{tabular}
\end{center}

\qn{2} A square loop is placed near a long straight current carrying wire as shown.
\qfig{p1_phy_s3_q02_1}
Match the following\par\nobreak
\begin{center}\small\setlength{\tabcolsep}{3pt}
\begin{tabular}{|c|p{2.5cm}|c|p{2.9cm}|}
\hline
\multicolumn{2}{|c|}{Column$-$I} & \multicolumn{2}{c|}{Column$-$II} \\
\hline
(A) & If current is increased & (P) & Induced current in loop is clockwise \\
\hline
(B) & If current is decreased & (Q) & Induced current in loop is anticlockwise \\
\hline
(C) & If loop is moved away from the wire & (R) & wire will attract the loop \\
\hline
(D) & If loop is moved toward the wire & (S) & wire will repel the loop \\
\hline
\end{tabular}
\end{center}

\qn{3}\par\nobreak
\begin{center}\small\setlength{\tabcolsep}{3pt}
\begin{tabular}{|c|p{2.2cm}|c|p{3.3cm}|}
\hline
\multicolumn{2}{|c|}{Column$-$I} & \multicolumn{2}{c|}{Column$-$II} \\
\hline
(A) & $\beta$- rays & (P) & Continuous energy distribution with dark lines in between \\
\hline
(B) & $\gamma$- rays & (Q) & Continuous energy distribution with sharp peaks in between \\
\hline
(C) & Absorption spectrum of Hydrogen & (R) & Continuous energy distribution \\
\hline
(D) & X-rays & (S) & Discrete energy distribution \\
\hline
\end{tabular}
\end{center}

\qn{4}\par\nobreak
\begin{center}\small\setlength{\tabcolsep}{3pt}
\begin{tabular}{|c|p{3.6cm}|c|p{2.0cm}|}
\hline
\multicolumn{2}{|c|}{Column$-$I} & \multicolumn{2}{c|}{Column$-$II} \\
\hline
(A) & Uniform positively charged Sphere having a spherical cavity. The electric\newline \qfigi{p1_phy_s3_q03_1}\newline (C$_1$ \& C$_2$ represents the center of sphere and spherical cavity respectively) & (P) & Field at C$_1$ is from C$_2$ to C$_1$ \\
\hline
(B) & Spherical cavity in a uniform spherical Planet. The gravitational\newline \qfigi{p1_phy_s3_q03_2}\newline (C$_1$ \& C$_2$ represents the center of sphere and spherical cavity respectively) & (Q) & Field at C$_1$ is from C$_1$ to C$_2$ \\
\hline
(C) & Cylindrical cavity in a wire carrying uniform current density. The magnetic\newline \qfigi{p1_phy_s3_q03_3}\newline (C$_1$ \& C$_2$ are the points on the axis of cylinder and cylindrical cavity respectively) & (R) & Field at C$_2$ is from C$_2$ to C$_1$ \\
\hline
(D) & Two uniformly charged non conducting shells outer shell is positively charged inner shell is negatively charged. The electric\newline \qfigi{p1_phy_s3_q03_4}\newline (C$_1$ \& C$_2$ represents the center of spherical shells) & (S) & Field at C$_2$ is perpendicular to C$_1$C$_2$ \\
\hline
\end{tabular}
\end{center}

\qn{5} A satellite is in a circular equatorial orbit of radius 7000 km around the Earth. If it is transferred to a circular orbit of double the radius,\par\nobreak
\begin{center}\small\setlength{\tabcolsep}{3pt}
\begin{tabular}{|c|p{2.8cm}|c|p{2.4cm}|}
\hline
\multicolumn{2}{|c|}{Column$-$I} & \multicolumn{2}{c|}{Column$-$II} \\
\hline
(A) & Angular momentum & (P) & Increases \\
\hline
(B) & Area of Earth covered by satellite signal & (Q) & Decreases \\
\hline
(C) & Potential energy & (R) & Magnitude becomes double \\
\hline
(D) & Kinetic energy & (S) & Magnitude becomes half \\
\hline
\end{tabular}
\end{center}

\partsec{2}{CHEMISTRY}
\sect{Section-I}

\qn{1} K$_{a_1}$ and K$_{a_2}$ are the 1$^{\text{st}}$ and 2$^{\text{nd}}$ ionization constants of maleic acid whereas K$'_{a_1}$ and K$'_{a_2}$ are those of fumaric acid, then choose the correct options among the following .
\begin{qoptions}
\opt{A}{$pK_{a_1}>pK'_{a_1}$}
\opt{B}{$pK_{a_2}>pK'_{a_2}$}
\opt{C}{$pK_{a_1}<pK'_{a_1}$}
\opt{D}{$pK_{a_2}<pK'_{a_2}$}
\end{qoptions}

\qn{2} In a body centred cubic unit cell of A$_2$ type (R is the radius of one particle at a lattice point)
\begin{qoptions}
\opt{A}{The edge length is equal to $4R\sqrt3$}
\opt{B}{In this arrangement face centre is a regular octahedral void}
\opt{C}{The edge length is equal to $\dfrac{2d}{\sqrt3}$ where `d' is the shortest distance between two lattice points}
\opt{D}{The square of the edge length is equal to $16R^2$}
\end{qoptions}

\qn{3} A disconnection of 1, 4-diphenylbut -2- yne is given below followed by possible corresponding synthetic reactions. Which of (A) $-$ (D) is/are include as suitable first reagent ?
\qfig{p1_chem_s1_q03_1}
Synthetic steps for the disconnect shown
\qfig{p1_chem_s1_q03_2}
\begin{qoptions}
\opt{A}{BuLi, THF}
\opt{B}{EtMgBr, Et$_2$O}
\opt{C}{NaOEt, EtOH}
\opt{D}{NaNH$_2$, NH$_3$}
\end{qoptions}

\qn{4} Which of the following statement(s) is/are true about the reaction given below?
\qfig{p1_chem_s1_q04_1}
\begin{qoptions}
\opt{A}{It involves a carbocation intermediate}
\opt{B}{Rearrangement is due to SN1 reaction mechanism}
\opt{C}{It proceeds via a concerted SN2 pathway}
\opt{D}{It involves neighbouring group participation}
\end{qoptions}

\qn{5} The potential energy of a diatomic molecule, as a function of the inter nuclear separation r, is approximated as $U(r)=\dfrac{A}{r^a}-\dfrac{B}{r^b}$ Where A \& B are positive constants and a $>$ b
\qfig{p1_chem_s1_q05_1}
As shown in the above figure r$^0$ is the equilibrium bond length. What is the energy necessary to break the bond from its equilibrium position
\begin{qoptions}
\opt{A}{$\dfrac{A}{r_0^a}-\dfrac{B}{r_0^b}$}
\opt{B}{$\dfrac{B}{r_0^b}-\dfrac{A}{r_0^a}$}
\opt{C}{$\dfrac{A}{r_0^a}\left(\dfrac ab-1\right)$}
\opt{D}{Invalid Data}
\end{qoptions}

\sect{Section-III}

\qn{1} At 25$^\circ$C, the degree of dissociation of pure water is $1.9\times10^{-9}$. If the specific conductance of water at this temperature is x $\times10^{-6}$ Sm$^{-1}$. Calculate the value of x. Given ionic molar conductances at infinite dilution for H$^+$ and OH$^-$ ions are $350\times10^{-4}$ and $200\times10^{-4}$ Sm$^2$ mole$^{-1}$ respectively

\qn{2} Identify the value of ( X $-$ Y )
\qfig{p1_chem_s2_q02_1}

\qn{3} Total no of products ( including stereoisomers ) are formed when (R) $-$ 3 $-$ methyl -1- pentene is treated with following
\qfig{p1_chem_s2_q03_1}
I ) Hg (OAc)$_2$, H$_2$O followed by NaBH$_4$ \quad II) H$_2$ / Pt\\
III) BH$_3$ followed by H$_2$O$_2$ in NaOH \quad IV) Br$_2$/CCl$_4$

\qn{4} Number of reactions that lead to the formation of ``pyro'' form of acid\\[2pt]
A) H$_3$PO$_4\ \xrightarrow[\text{510K}]{\Delta}$\\[2pt]
B) PCl$_3$ + H$_3$PO$_3\to$\\[2pt]
C) Phosphorous acid + Br$_2\ \xrightarrow[\text{tube}]{\text{Sealed}}$ A $\xrightarrow[\text{H}_2\text{O}]{\text{Excess}}$ B $\xrightarrow{\Delta}$ C\\[2pt]
D) P$_4$O$_{10}$ + PCl$_3\to$ D $\xrightarrow[\text{excess}]{\text{H}_2\text{O}}$ E $\xrightarrow{\Delta}$ F\\[2pt]
E) 3RCOOH + PCl$_3\to$ 3RCOCl + [w]\\
\hspace*{1.5em}[w] $\xrightarrow{\text{H}_2\text{SO}_4}$ Y$_g\uparrow$ +X + H$_2$O\\
\hspace*{1.5em}[x] $\xrightarrow{\Delta}$ Z

\qn{5} Number of reactions that result formation of precipitate [except white ],finally.\\[2pt]
A) Pb(NO$_3$)$_2\ \xrightarrow{\text{NaOH}}$ A $\xrightarrow{\text{NaOH}}$ B $\xrightarrow{\text{H}_2\text{S}}$ C\\[2pt]
B) FeCr$_2$O$_4\ \xrightarrow{\text{O}_2+\text{Na}_2\text{CO}_3}$ $\xrightarrow{\text{reaction}}$ solution $\xrightarrow[\text{solution}]{\text{AgNO}_3}\xrightarrow{\text{Ammonia}}$\\[2pt]
C) Fe$^{+3}$ + K$_3$[Fe(CN)$_6$] $\to$ D $\xrightarrow{\text{H}_2\text{O}_2}$ E\\[2pt]
D) Zn$^{+2}$ + K$_4$[Fe(CN)$_6$] $\to$ F $\xrightarrow[\text{acid}]{\text{dilute}}$
\qfig{p1_chem_s2_q05_1}
\qfig{p1_chem_s2_q05_2}

\sect{Section-IV}

\qn{1} Match column I with correct thermodynamic parameters in Column $-$ II (symbols carry the usual meaning, assume gases to be ideal).\par\nobreak
\begin{center}\small\setlength{\tabcolsep}{3pt}
\begin{tabular}{|c|p{3.9cm}|c|p{1.5cm}|}
\hline
\multicolumn{2}{|c|}{Column$-$I} & \multicolumn{2}{c|}{Column$-$II} \\
\hline
(A) & One mole of nitrogen gas contained in a cylinder at 300 K is allowed to expand isothermally against an external pressure of 5 atm. & (P) & w $>$ 0 \\
\hline
(B) & 3 moles of helium gas at 1 atm are compressed reversibly and isothermally at 400 K to 5 atm pressure & (Q) & $\Delta$U = 0 \\
\hline
(C) & An ideal gas is compressed reversibly and adiabatically from a volume of 1dm$^3$ to a volume of 0.25 dm$^3$ & (R) & $\Delta$H = 0 \\
\hline
(D) & One mole of liquid benzene is converted reversibly into vapour at its boiling point by heating it. & (S) & $\Delta$H $>$ $\Delta$U \\
\hline
\end{tabular}
\end{center}

\qn{2} Column $-$ I has unbalanced reaction, Column $-$ II has the relation between equivalent mass (E) and molecular mass (M) of the underlined species in column $-$ I.\par\nobreak
\begin{center}\small\setlength{\tabcolsep}{3pt}
\begin{tabular}{|c|p{3.3cm}|c|p{2.1cm}|}
\hline
\multicolumn{2}{|c|}{\parbox{3.6cm}{\centering Column$-$I\\Unbalanced Reaction}} & \multicolumn{2}{c|}{Column$-$II} \\
\hline
(A) & $\underline{\text{P}_2\text{H}_4}\to$ PH$_3$ + P$_4$H$_2$ & (P) & $E=\dfrac{3M}{4}$ \\
\hline
(B) & $\underline{\text{I}_2}$ + HNO$_3\to$ HIO$_3$ + NO$_2$ + H$_2$O & (Q) & $E=\dfrac M3$ \\
\hline
(C) & $\underline{\text{CuO}}$ + $\underline{\text{NH}_3}\to$ N$_2$ + Cu + H$_2$O & (R) & $E=\dfrac{5M}{6}$ \\
\hline
(D) & H$_2$SO$_3$ + H$_2$S $\to$ $\underline{\text{S}}$ + H$_2$O & (S) & stoichiometric coefficient of one of the reactant species is 1 \\
\hline
\end{tabular}
\end{center}

\qn{3} Match the compound in Column $-$I with their characteristic tests(s)/reaction(s) given in Column $-$II\par\nobreak
\begin{center}\small\setlength{\tabcolsep}{3pt}
\begin{tabular}{|c|p{3.3cm}|c|p{2.5cm}|}
\hline
\multicolumn{2}{|c|}{Column$-$I} & \multicolumn{2}{c|}{Column$-$II} \\
\hline
(A) & H$_2$N $-$ $\overset{\oplus}{\text{N}}$H$_3\overset{\ominus}{\text{Cl}}$ & (P) & Sodium fusion extract of the compound gives Prussian blue colour with FeSO$_4$ \\
\hline
(B) & \qfigi{p1_chem_s3_q03_B} & (Q) & gives positive FeCl$_3$ test \\
\hline
(C) & \qfigi{p1_chem_s3_q03_C} & (R) & gives white precipitate with AgNO$_3$ \\
\hline
(D) & \qfigi{p1_chem_s3_q03_D} & (S) & reacts with aldehydes to form the corresponding hydrazone derivative \\
\hline
\end{tabular}
\end{center}

\qn{4} Match column$-$I with correct thermodynamic parameters in Column $-$ II (symbols carry the usual meaning, assume gases to be ideal).\par\nobreak
\begin{center}\small\setlength{\tabcolsep}{3pt}
\begin{tabular}{|c|p{3.2cm}|c|p{2.8cm}|}
\hline
\multicolumn{2}{|c|}{Column$-$I} & \multicolumn{2}{c|}{Column$-$II} \\
\hline
(A) & [Co(NH$_3$)$_4$ (H$_2$O)$_2$] Cl$_2$ & (P) & Geometrical isomers \\
\hline
(B) & [Pt (NH$_3$)$_2$ Cl$_2$] & (Q) & Paramagnetic \\
\hline
(C) & [Co(H$_2$O)$_5$ Cl] Cl & (R) & Diamagnetic \\
\hline
(D) & [Ni(H$_2$O)$_6$] Cl$_2$ & (S) & On reaction with AgNO$_3$ forms 2 equivalents of AgCl \\
\hline
\end{tabular}
\end{center}

\qn{5} Match the following:\par\nobreak
\begin{center}\small\setlength{\tabcolsep}{3pt}
\begin{tabular}{|c|p{1.5cm}|c|p{4.6cm}|}
\hline
\multicolumn{2}{|c|}{Column$-$I} & \multicolumn{2}{c|}{Column$-$II} \\
\hline
(A) & XeF$_2$ & (P) & Undergoes hydrolysis with water \\
\hline
(B) & XeF$_4$ & (Q) & Act as oxidizing agent \\
\hline
(C) & XeF$_6$ & (R) & Undergoes addition reaction \\
\hline
(D) & XeO$_3$ & (S) & Gives disproportionation reaction with water or hydroxide ion. \\
\hline
\end{tabular}
\end{center}

\partsec{3}{MATHEMATICS}
\sect{Section-I}

\qn{1} F, g : $\mathbb R/\{1\}\to\mathbb R$ such that $f(2x+1)+2g(2x+1)=2x$ and $f\left(\dfrac{x}{x-1}\right)+h\left(\dfrac{x}{x-1}\right)=x$, then
\begin{qoptions}
\opt{A}{$f(x)=\dfrac{x^2-4x+1}{x-1}$}
\opt{B}{$g(x)=\dfrac{x^2-3x+1}{x-1}$}
\opt{C}{$f(5)=\dfrac{-3}{2}$}
\opt{D}{$g(2)=1$}
\end{qoptions}

\qn{2} Let A, B $\in M_n(\mathbb R)$ are square matrices of same order such that AB $-$ BA = A. Then
\begin{qoptions}
\opt{A}{$A^5B-BA^5=5A^5$}
\opt{B}{trace (A) = 0}
\opt{C}{A must be an Idempotent Matrix}
\opt{D}{A must be Nilpotent Matrix}
\end{qoptions}

\qn{3} Let f: [0, 1] $\to\mathbb R$ be a continuous function such that for all non-negative integer n, $\displaystyle\int_0^1x^nf(x)\,dx=0$. Then
\begin{qoptions}
\opt{A}{f (0) = 0}
\opt{B}{f (x) must be identically zero}
\opt{C}{f(1) = 1}
\opt{D}{$\exists\,x_1,x_2\in[0,1]$ such that $f(x_1)f(x_2)<0$}
\end{qoptions}

\qn{4} Which of the following is/are true.
\begin{qoptions}
\opt{A}{f, g : $\mathbb R\to\mathbb R$, $g(x)\ge0$, f(x) is twice differentiable and $f''(x)+f(x)=-xg(x)f'(x)\Rightarrow f(x)$ is a bounded function}
\opt{B}{f: $\mathbb R\to[-1,1]$, twice differentiable such that $(f(0))^2+(f'(0))^2=4$, then $\exists\ x_0\in R$ such that $f(x_0)+f''(x_0)=0$}
\opt{C}{Number of solutions of an equation $2^x=1+x^2$ is two.}
\opt{D}{None of these}
\end{qoptions}

\qn{5} Let two circles of radius r$_1$ and r$_2$ (where r$_1\ne$ r$_2$) intersect each other at point P and Q. Consider the direct common tangent of the two circles, nearer to point P, touching them at points A and B respectively then
\begin{qoptions}
\opt{A}{Circum radius of $\Delta$PAB $=\sqrt{r_1r_2}$}
\opt{B}{Circum radius of $\Delta$QAB $=\sqrt{r_1r_2}$}
\opt{C}{Circum radius of $\Delta$QAB $=\dfrac{r_1+r_2}{2}$}
\opt{D}{Circum radius of $\Delta$PAB $<$ circum radius of $\Delta$QAB}
\end{qoptions}

\sect{Section-III}

\qn{1} If a function f(x) satisfies $f(x+1)=4f(x)(1-f(x))$ and number of possible choices for f(1) such that f(2017) = 0, is K. Then sum of all digits of $\lfloor\log_2K\rfloor$ is \underline{\hspace{1.5cm}}

\qn{2} In the largest possible side of an equilateral triangle with vertices within a unit square (the vertices of the triangle can be inside the square or on its boundary) is $\sqrt a-\sqrt b$; $a,b\in\mathbb N$. Then a + b is ....

\qn{3} If $\displaystyle\lim_{n\to\infty}\left(((n+1)!)^{\frac1{n+1}}-(n!)^{\frac1n}\right)=\dfrac ab$, then $\left\lfloor\dfrac ba\right\rfloor$ is ............

\qn{4} Let f : $\mathbb R\to\mathbb R$ be twice continuously differentiable function such that $2f(x+1)=f(x)+f(2x)$, then $\left\lfloor\dfrac{f(2)+f(7)}{f(9)}\right\rfloor$ is \underline{\hspace{1.5cm}}

\qn{5} f : $\mathbb R\to\mathbb R$ be a continuous function such that $f(f(x))+f(x)+x=0\ \forall\ x\in\mathbb R$ then number of such function is ..........

\sect{Section-IV}

\qn{1}\par\nobreak
\begin{center}\small\setlength{\tabcolsep}{3pt}
\begin{tabular}{|c|p{4.4cm}|c|p{0.5cm}|}
\hline
\multicolumn{2}{|c|}{Column-I} & \multicolumn{2}{c|}{Column-II} \\
\hline
(A) & The value of $\left|\dfrac{\tan40^\circ+\tan100^\circ+\tan160^\circ}{\tan20^\circ\tan40^\circ\tan80^\circ}\right|$ is & (P) & 6 \\
\hline
(B) & If $\alpha,\beta$ are roots of $ax^2+bx+c=0$, $a\ne0$ and $\dfrac\alpha\beta,\dfrac\beta\alpha$ are roots of $x^2-3x+1=0$, then the value of $\dfrac{b^2}{ac}$ is & (Q) & 5 \\
\hline
(C) & If P is the product of solutions of the equation $x^{\ell nx}=e^2x$, then P is less than & (R) & 3 \\
\hline
(D) & If the expression $kx^2+5xy+2y^2+3x+3y+1$ can be split into two linear factors, then the value of $|k|$ is greater than or equal to & (S) & 2 \\
\hline
\end{tabular}
\end{center}

\qn{2}\par\nobreak
\begin{center}\small\setlength{\tabcolsep}{3pt}
\begin{tabular}{|c|p{3.3cm}|c|p{2.4cm}|}
\hline
\multicolumn{2}{|c|}{Column-I} & \multicolumn{2}{c|}{Column-II} \\
\hline
(A) & The set of points z satisfying $|z-i||z||=|z+i||z||$, is contained in or equal to & (P) & An ellipse with eccentricity $\dfrac45$ \\
\hline
(B) & The set of points z satisfying $|z+4|+|z-4|=10$, is contained in or equal to & (Q) & The set of points z satisfying Im z = 0 \\
\hline
(C) & If $|\omega|=2$, then the set of points $z=\omega-\dfrac1\omega$ is contained in or equal to & (R) & The set of points z satisfying Imz $\le1$ \\
\hline
(D) & If $|\omega|=1$, then the set of points $z=\omega+\dfrac1\omega$ is contained in or equal to & (S) & The set of points z satisfying $|\text{Re}z|\le2$ \\
\hline
\end{tabular}
\end{center}

\qn{3} Match the following for system of linear equations:-\\
$2x-3y+5z=12$, $3x+y+\lambda z=\mu$, $x-7y+8z=17$\par\nobreak
\begin{center}\small\setlength{\tabcolsep}{3pt}
\begin{tabular}{|c|p{3.3cm}|c|p{2.3cm}|}
\hline
\multicolumn{2}{|c|}{Column-I} & \multicolumn{2}{c|}{Column-II} \\
\hline
(A) & Unique solution & (P) & $\lambda=2,\ \mu=7$ \\
\hline
(B) & Infinite solutions & (Q) & $\lambda\ne2,\ \mu=7$ \\
\hline
(C) & No solution & (R) & $\lambda\ne2,\ \mu\ne7$ \\
\hline
(D) & Consistent system of equation & (S) & $\lambda=2,\ \mu\ne7$ \\
\hline
\end{tabular}
\end{center}

\qn{4} Let $\hat a,\hat b$ and $\hat c$ are three unit vectors mutually inclined to each other at an angle $\theta$. Also $\hat d$ is an unit vector such that $\hat d\cdot\hat a=\hat d\cdot\hat b=\hat d\cdot\hat c=\cos\alpha$.\\
Match the entries for $\theta$ given in Column-I with their corresponding values for $\alpha$ given in Column-II.\par\nobreak
\begin{center}\small\setlength{\tabcolsep}{3pt}
\begin{tabular}{|c|c|c|c|}
\hline
\multicolumn{2}{|c|}{Column-I ($\theta$)} & \multicolumn{2}{c|}{Column-II ($\alpha$)} \\
\hline
(A) & $\dfrac\pi2$ & (P) & $\dfrac\pi4$ \\
\hline
(B) & $\dfrac\pi3$ & (Q) & $\cos^{-1}\left(\dfrac1{\sqrt3}\right)$ \\
\hline
(C) & $\cos^{-1}\left(\dfrac14\right)$ & (R) & $\cos^{-1}\sqrt{\dfrac23}$ \\
\hline
(D) & $\cos^{-1}\left(\dfrac55\right)$ & (S) & $\dfrac\pi6$ \\
\hline
\end{tabular}
\end{center}

\qn{5}\par\nobreak
\begin{center}\small\setlength{\tabcolsep}{3pt}
\begin{tabular}{|c|p{4.6cm}|c|p{0.8cm}|}
\hline
\multicolumn{2}{|c|}{Column-I} & \multicolumn{2}{c|}{Column-II} \\
\hline
(A) & Number of ways in which 3 squares of unit length can be chosen on a $8\times8$ chessboard, so that all squares are in same diagonal line, is & (P) & 441 \\
\hline
(B) & Number of three letter words that can be formed using the letters of the word ``LIVELIHOOD'', is & (Q) & 361 \\
\hline
(C) & Number of five element subset that can be chosen from the set of the first 11 natural numbers so that atleast two of the five elements are consecutive, is & (R) & 392 \\
\hline
(D) & Number of ways in which the letters of the word `ABBCABBC' can be arranged such that the sequence of the letters ABBC does not appear in any word, is & (S) & 264 \\
\hline
\end{tabular}
\end{center}


\papersec{2}
\partsec{1}{PHYSICS}
\sect{Section-I(i)}

\qn{1} A large wooden plate of area $10\ \text{m}^2$ floating on the surface of a river is made to move horizontally with a speed of 2 m/s by applying a tangential force. River is 1 m deep and the water in contact with the bed is stationary. Then choose correct statement(s). (Coefficient of viscosity of water $=10^{-3}$ N-s/m)
\begin{qoptions}
\opt{A}{velocity gradient is $2\,\text{s}^{-1}$}
\opt{B}{velocity gradient is $1\,\text{s}^{-1}$}
\opt{C}{force required to keep the plate moving with constant speed is 0.02 N}
\opt{D}{force required to keep the plate moving with constant speed is 0.01 N}
\end{qoptions}

\qn{2} A ladder AB is supported by a smooth vertical wall and sufficient rough horizontal floor as shown. A boy starts moving from A to B slowly. The ladder remains at rest, the pick up the correct statement(s):
\qfig{p2_phy_s1_q02_1}
\begin{qoptions}
\opt{A}{Magnitude of normal reaction by wall on ladder at point B will increase.}
\opt{B}{Magnitude of normal reaction by wall on ladder at point A will decrease}
\opt{C}{Magnitude of normal reaction by floor on ladder at point A will remain unchanged.}
\opt{D}{Magnitude of friction force by floor on ladder at point A will increase}
\end{qoptions}

\qn{3} In which of the following cases will an erect image be formed
\begin{qoptions}
\opt{A}{An object between pole and focus of a concave mirror}
\opt{B}{An object between pole and object focus of a converging lens}
\opt{C}{An object in front of a plane mirror}
\opt{D}{An object between object focus of a converging lens and infinity}
\end{qoptions}

\qn{4} An orbital electron in a particular state of hydrogen has the radius R. This orbital electron is excited to another state by some energy transfer to the hydrogen atom so that its radius is now 9R. Then
\begin{qoptions}
\opt{A}{Its energy is now one ninth of the original energy}
\opt{B}{Its velocity is now one third of the original velocity}
\opt{C}{Its angular momentum is three times the original angular momentum}
\opt{D}{It will come to the original state by emitting maximum of 2 photons}
\end{qoptions}

\qn{5} A straight conducting bar AB of mass m and length $l$ is suspended horizontal by two identical non-conducting springs of spring constant k as shown in the figure. The bar is located in a region of uniform magnetic field of intensity B. The capacitor C is initially charged to the potential V. Springs are in equilibrium position with zero potential energy. At time t = 0, the switch S is closed and capacitor discharge. The bar starts oscillating in the vertical plane. (Assume that the time of discharge of the capacitor is much smaller than the time period of mechanical oscillation of the bar. Neglect the effect of gravity). Choose the correct statement(s).
\qfig{p2_phy_s1_q05_1}
\begin{qoptions}
\opt{A}{The maximum speed of bar during oscillation is $\dfrac{BlCV}{3m}$}
\opt{B}{The maximum speed of bar during oscillation is $\dfrac{BlCV}{m}$}
\opt{C}{The amplitude of mechanical oscillation of bar is $\sqrt{\dfrac{B^2l^2C^2V^2}{2mk}}$}
\opt{D}{The amplitude of mechanical oscillation of bar is $\sqrt{\dfrac{B^2l^2C^2V^2}{mk}}$}
\end{qoptions}

\qn{6} In three experiments, a material A at a particular low temperature $T_C$ and a material B at a particular high temperature $T_H$ are placed in an isolated and insulated container. When they reach thermal equilibrium with each other (No phase change occurs), their final temperature $T_f$ is measured. The masses $m_A$ \& $m_B$ and specific heats $C_A$ \& $C_B$ of the materials are given in the table. Assume that heat transferred is Q in the experiment.\par\nobreak
\begin{center}\small
\begin{tabular}{|c|c|c|c|c|}
\hline
Experiment & $m_A$ & $C_A$ & $m_B$ & $C_B$ \\
\hline
1 & m & c & m & c \\
\hline
2 & m & c & 2m & c \\
\hline
3 & m & 2c & m & c \\
\hline
\end{tabular}
\end{center}
Then which of the following is correct
\begin{qoptions}
\opt{A}{$(T_f)_1>(T_f)_2>(T_f)_3$}
\opt{B}{$Q_2>Q_1>Q_3$}
\opt{C}{$Q_1>Q_2>Q_3$}
\opt{D}{$(T_f)_2>(T_f)_1>(T_f)_3$}
\end{qoptions}

\sect{Section-I(ii)}

\paraq{7}{8}
In a gas where rms thermal speed is $v_t\left(\sqrt{\dfrac{3kT}{m}}\right)$, coulomb force will ensure that two nuclei each with charge q, cannot approach each other any closer than the distance of closest approach $r_m$. The value of $r_m$ is the distance at which their kinetic energy is completely converted to potential energy. Separation between two hydrogen nuclei in core of the sun cannot be small enough to allow fusion to occur. Known temperature that exists in the sun's core, coulomb force will cause the nuclei to remain at distance from each other equal to 100-1000 (or more) of nuclear radii $\left(\simeq10^{-15}\text{m}\right)$. So according to classical physics proton fusion reaction should not occur in sun. But if we consider quantum nature of nuclei, quantum-mechanical tunneling allows the protons to go under the coulomb barrier with a good probability. Bridging criteria for fusion was found to be $\lambda=\dfrac{h}{mv_t}\leqslant r_m$ where $\lambda$ is the de Broglie wavelength of moving proton when it is at large distance from the other proton. Wherever this equation is satisfied, two nuclei can approach each other arbitrarily closely so that nuclear force can lead to fusion. We find that the rms thermal speed in core of star must exceed a critical value for fusion to occur.

\qn{7} According to classical physics temperature(s) of sun core which allows the fusion.
\begin{qoptions}
\opt{A}{$10^{10}$ K}
\opt{B}{$10^{11}$ K}
\opt{C}{$10^{7}$ K}
\opt{D}{$10^{6}$ K}
\end{qoptions}

\qn{8} On considering quantum effects the core temperature(s) of sun which allows the fusion.
\begin{qoptions}
\opt{A}{$10^{4}$ K}
\opt{B}{$10^{5}$ K}
\opt{C}{$10^{8}$ K}
\opt{D}{$10^{7}$ K}
\end{qoptions}

\paraq{9}{10}
A student constructs a series RLC circuit. While operating the circuit with an ideal source of 60 Volt emf at a frequency $f_0$. She uses an ideal A.C. voltmeter and ideal A.C. ammeter to measure the potential difference across each device and current in circuit. The potential difference across resistor, inductor coil and capacitor is obtained as 24.00, 37.94V and 84.00V respectively. Net impedance of circuit at this frequency is $50\ \Omega$.

\qn{9} Select correct option(s)
\begin{qoptions}
\opt{A}{Reading of ammeter is 1.2 A}
\opt{B}{Reading of ammeter is 1.5 A}
\opt{C}{The combined potential difference across inductor coil and capacitor will be 49.47V}
\opt{D}{The combined potential difference across inductor coil and capacitor will be 51.52V}
\end{qoptions}

\qn{10} Select the correct option(s)
\begin{qoptions}
\opt{A}{Power factor of circuit is 0.8}
\opt{B}{Power factor of circuit is 0.6}
\opt{C}{resonance frequency of RLC circuit is $\sqrt{\dfrac73}f_0$}
\opt{D}{resonance frequency of RLC circuit is $\sqrt{\dfrac53}f_0$}
\end{qoptions}

\paraq{11}{12}
In a hydrogen atom, the electron revolves around the nucleus (proton) in circular path. In addition, then electron also spins about it's own axis. Due to this spin, electron has a magnetic dipole moment $(\vec\mu)$ in direction opposite to it's spin angular momentum $\left(\vec S\right)$. As seen from the frame of electron, the electron itself is only spinning, but the proton is moving around it in a circular path with the same speed and radius as electron in ground frame. This motion of proton around electron produces a magnetic field $\left(\vec B\right)$. This magnetic field and magnetic dipole moment of electron have a mutual potential energy. If orbital angular momentum $\left(\vec L\right)$ of electron and it's spin angular momentum $\left(\vec S\right)$ are parallel, the potential energy is different from the case when they are anti parallel.

\qn{11} Choose the correct statement(s)
\begin{qoptions}
\opt{A}{If $\vec L$ and $\vec S$ are parallel, $\vec\mu$ is parallel to $\vec B$}
\opt{B}{If $\vec L$ and $\vec S$ are parallel, $\vec\mu$ is anti parallel $\vec B$}
\opt{C}{If $\vec L$ and $\vec S$ are anti parallel, $\vec\mu$ is parallel to $\vec B$}
\opt{D}{If $\vec L$ and $\vec S$ are anti parallel, $\vec\mu$ is anti parallel to $\vec B$}
\end{qoptions}

\qn{12} Assuming Bohr's theory of hydrogen atom is valid. Select the correct option(s).
\begin{qoptions}
\opt{A}{the potential energy difference referred to in paragraph is inversely proportional to $n^3$}
\opt{B}{the potential energy difference referred to in paragraph is inversely proportional to $n^5$}
\opt{C}{the potential energy of parallel $\vec L$ and $\vec S$ is more than potential energy of anti parallel $\vec L$ and $\vec S$ case}
\opt{D}{the potential energy of parallel $\vec L$ and $\vec S$ case is less than potential energy of anti parallel $\vec L$ and $\vec S$ case}
\end{qoptions}

\sect{Section-III}

\qn{1} A solid right circular cone has its base cut out in conical shape shown in figure such that the hollow is a right circular cone on the same base. C.M. of the object coincide with the vertex of the hollow cone. Height of solid cone/height of hollow cone will be
\qfig{p2_phy_s3_q01_1}

\qn{2} A satellite is put in an orbit just above the earth's surface with a velocity $\sqrt{1.5}$ time the velocity for a circular orbit. The maximum distance of the satellite from the center of earth is nR where R is radius of earth. Find n

\qn{3} In a uniform gravitational field g, which is acting vertically downward, a ball is thrown at an angle with horizontal such that the initial (immediately after projection) radius of curvature is 8 times the minimum radius of curvature. If angle of projection is $10x$ degree, find the value of x.

\qn{4} What should be the value of resistance R (in $\Omega$) in the circuit shown in figure so that the electric body consumes the rated power?
\qfig{p2_phy_s3_q04_1}

\qn{5} When an object is placed at a distance of 25 cm from a mirror, the magnification of real image formed is $m_1$. The object is moved 15 cm away from the mirror with respect to the earlier position along principal axis, magnification of real image formed becomes $m_2$. If $m_1\times m_2=4$, find the power of mirror in diopter.

\qn{6} Two identical hoop of same mass is rigidly attached at their circumference as shown. Now the system is hinged from point `P' such that it can oscillate. Let $T_1$ be the period of oscillation, if system is slightly displaced from its equilibrium position and oscillation occurs in xy plane. Now system is turned through $90^\circ$ such that plane of hoop is yz plane. Let $T_2$ be the period of oscillation, if system is again slightly displaced from its equilibrium position and oscillation occurs again in xy plane. If the ratio of period of oscillation $\left(\dfrac{T_1}{T_2}\right)$ is $\sqrt{\dfrac{36}{11\lambda}}$. Then find the value of $\lambda$.
\qfig{p2_phy_s3_q06_1}

\qn{7} Two infinitely long wires are placed perpendicular to the plane of paper. Current in wire A is $4i_0$ outward the plane of paper and in `B' is $i_0$ inward the plane of paper. The $\displaystyle\int_{-\infty}^{+\infty}\vec B\cdot d\vec l$ along the QP is $\dfrac K2\mu_0 i_0$. Find the value of K.
\qfig{p2_phy_s3_q07_1}

\qn{8} Two charge particles are fixed at $P_1$ and $P_2$ separated by a distance 9 m. Find position (in meter) on the line $P_1P_2$ from $P_1$ where electric field is zero (null point).
\qfig{p2_phy_s3_q08_1}

\qn{9} Intensity observed in an interference pattern is $I=I_0\sin^2\theta$. At $\theta=30^\circ$, intensity $I=5\pm0.0020\ \text{W/m}^2$. If $I_0=20\ \text{W/m}^2$ and percentage error in angle $\theta$ is $\dfrac n\pi\sqrt3\times10^{-2}\%$, then the value of n is.

\qn{10} A uniform disc of radius r and mass m is charged uniformly with the charge q. This disc is placed flat on a rough horizontal surface having coefficient of friction $\mu$. A uniform magnetic field is present in a circular region $(a>r)$ but varying as $kt^3$ as shown in figure. Find the time in second after which the disc begins to rotate.\\
(Given r = 1 m, m = 18 kg, q = 1C, $\mu=0.1$m, K = 4, g = 10 m/s$^2$)
\qfig{p2_phy_s3_q10_1}

\partsec{2}{CHEMISTRY}
\sect{Section-I(i)}

\qn{1} $V_1$ ml of 3M HCl solution is mixed with $V_2$ ml of 3M H$_2$SO$_4$ to form a solution in which molarity of H$^+$ ions is 5 M, then the value of $V_1$ and $V_2$ may be
\begin{qoptions}
\opt{A}{150 ml, 75 ml}
\opt{B}{150 ml, 300 ml}
\opt{C}{125 ml, 250 ml}
\opt{D}{100 ml, 200 ml}
\end{qoptions}

\qn{2} At high temperature, gaseous compound S$_4$N$_4$ decomposes into nitrogen gas and sulphur vapours. It was found that 20 ml of gaseous compound S$_4$N$_4$ decomposed forming 50 ml of total gaseous products. Based on the above information, which of the following options is/are correct?
\begin{qoptions}
\opt{A}{Volume of nitrogen gas formed is 30 ml}
\opt{B}{Volume of sulphur vapours formed is 20 ml}
\opt{C}{formula of sulphur vapours formed is S$_8$}
\opt{D}{The volume of sulphur vapours obtained at any time would be twice the vapours of S$_4$N$_4$ decomposed.}
\end{qoptions}

\qn{3} Which is/are the resonance structure of methyl isocyanate ?
\begin{qoptions}
\opt{A}{\qfigi{p2_chem_s1_q03_A}}
\opt{B}{\qfigi{p2_chem_s1_q03_B}}
\opt{C}{\qfigi{p2_chem_s1_q03_C}}
\opt{D}{\qfigi{p2_chem_s1_q03_D}}
\end{qoptions}

\qn{4} Which of the following product(s) is(are) formed when 1,2,4-trimethylcyclopentene is hydrogenated?
\begin{qoptions}
\opt{A}{\qfigi{p2_chem_s1_q04_A}}
\opt{B}{\qfigi{p2_chem_s1_q04_B}}
\opt{C}{\qfigi{p2_chem_s1_q04_C}}
\opt{D}{\qfigi{p2_chem_s1_q04_D}}
\end{qoptions}

\qn{5} $\text{P}\xrightarrow{\text{aq.NO}_2^-}\underset{\text{brown}}{\text{Q}}\xrightarrow{\text{aq.S}^{-2}}\underset{\text{purple}}{\text{R}}$\\
Choose the correct statement(s)
\begin{qoptions}
\opt{A}{Number of unpaired electrons in [P] is zero}
\opt{B}{Oxidation state of central metal atom / ion in [Q] is [+3]}
\opt{C}{Colour of R arises due to charge transfer}
\opt{D}{P, Q are capable to show only structural isomerism but R can show stereo isomerism also}
\end{qoptions}

\qn{6} Which statement(s) is/are true for BX$_3$ [X = F, Cl, Br, I]
\begin{qoptions}
\opt{A}{BF$_3$ reacts faster with lewis base Cl$^-$ than BCl$_3$}
\opt{B}{Rate of hydrolysis of BF$_3$ increases as F is replaced by H$_2$O}
\opt{C}{$\pi$-charge transfer from X to B to increase in the sequence F $<$ Cl $<$ Br $<$ I.}
\opt{D}{$\pi$-bond energy in BX$_3$ follows the order BF$_3>$ BCl$_3>$ BBr$_3>$ BI$_3$}
\end{qoptions}

\sect{Section-I(ii)}

\paraq{7}{8}
Colloids are thermodynamically unstable as it involves creation of surface. Their stability can be attributed to various kinetic factors, like charge etc. Solvation is also another factor that stabilizes lyophilic sols.

\qn{7} A negatively charged sol is formed when
\begin{qoptions}
\opt{A}{AgNO$_3$ is treated with excess of KI}
\opt{B}{SnO$_2$ is treated with NaOH}
\opt{C}{H$_2$S is passed into As$_2$O$_3$}
\opt{D}{FeCl$_3$ is added to a fresh precipitate of Fe(OH)$_3$}
\end{qoptions}

\qn{8} Which of the following are related to lyophobic sols
\begin{qoptions}
\opt{A}{Viscosity is same as that of the medium}
\opt{B}{Extensive hydration takes place}
\opt{C}{Migration of particles under the influence of electric field}
\opt{D}{Very difficult to detect particles even under ultra-microscope}
\end{qoptions}

\paraq{9}{10}
An acyclic hydrocarbon P, having molecular formula C$_6$H$_{10}$, gave acetone as the only organic product through the following sequence of reaction, in which Q is an intermediate organic compound.
\qfig{p2_chem_s2_para9_10_1}

\qn{9} The structure of compound P is
\begin{qoptions}
\opt{A}{CH$_3$CH$_2$CH$_2$CH$_2$ $-$ C $\equiv$ C $-$ H}
\opt{B}{H$_3$CCH$_2$ $-$ C $\equiv$ C $-$ CH$_2$CH$_3$}
\opt{C}{\qfigi{p2_chem_s2_q09_C}}
\opt{D}{\qfigi{p2_chem_s2_q09_D}}
\end{qoptions}

\qn{10} The structure of compound Q is
\begin{qoptions}
\opt{A}{\qfigi{p2_chem_s2_q10_A}}
\opt{B}{\qfigi{p2_chem_s2_q10_B}}
\opt{C}{\qfigi{p2_chem_s2_q10_C}}
\opt{D}{\qfigi{p2_chem_s2_q10_D}}
\end{qoptions}

\paraq{11}{12}
Halogens form many oxides with oxygen but most of them are unstable. Fluorine forms two compounds OF$_2$ and O$_2$F$_2$. Chlorine, bromine \& Iodine form oxides in which the oxidation states of these halogens range from +1 to +7. Chlorine oxides are highly reactive oxidising agents.

\qn{11} Choose the incorrect statements from the following $-$
\begin{qoptions}
\opt{A}{ClO$_2$ has more dimerising tendency than NO$_2$.}
\opt{B}{In Cl$_2$O$_7$ Skelton\\ \qfigi{p2_chem_s2_q11_1}\\ O is sp$^2$ hybridized \& $\theta_1>\theta_2$, x $>$ y}
\opt{C}{I$_2$O$_4$ on heating decompose to give I$_2$}
\opt{D}{OF$_2$ is called fluorine oxide.}
\end{qoptions}

\qn{12} Which of the following undergoes disproportionation in H$_2$O ?
\begin{qoptions}
\opt{A}{F$_2$}
\opt{B}{Cl$_2$}
\opt{C}{Br$_2$}
\opt{D}{I$_2$}
\end{qoptions}

\sect{Section-III}

\qn{1} The de-Broglie wavelength of the electron in a certain Bohr's orbit of hydrogen atom is 9.985 \AA. What is the principal quantum number (n) of this orbit?

\qn{2} The following graph shows two isotherms for a fixed mass of an ideal gas (x). The r.m.s speed of molecules of gas (x) at temperature $T_1$ and $T_2$ are $C_1$ and $C_2$ respectively. The value of $C_2/C_1$ is \underline{\hspace{1cm}}
\qfig{p2_chem_s3_q02_1}

\qn{3} A 0.001 molal solution of complex Pt(NH$_3$)$_4$Cl$_4$ in water had a freezing point depression $0.0056^\circ$C. $K_f$ for water is 1.86 K kg mole$^{-1}$. The number of moles of AgNO$_3$ required to react with one mole of the complex in aqueous solution is

\qn{4} Find the total number of possible monobromo products obtained during the reaction
\qfig{p2_chem_s3_q04_1}

\qn{5} The amount (in moles) of NaOH required for complete neutralization of one mole of the following compound is \underline{\hspace{1cm}}
\qfig{p2_chem_s3_q05_1}

\qn{6} How many structures given below from a $-$ e that cannot be a reactant X of the following ozonolysis sequence?
\qfig{p2_chem_s3_q06_1}
\qfig{p2_chem_s3_q06_2}

\qn{7} In the given reaction sequence, all possible five and six member lactones are X. Find X $-$ 5
\begin{center}
`X' (lactone), molecular formula C$_6$H$_{10}$O$_2$ $\xrightarrow[\text{(ii) H}_2\text{O};\ \text{(iii) Red P/HI}]{\text{(i) LiAlH}_4}$ \qfigi{p2_chem_s3_q07_1}
\end{center}

\qn{8} How many of the following statements are correct ?\\
A) Mg$^{+2}$ ion form more stable complex with EDTA as compare to Ca$^{+2}$\\
B) C$_{60}$ contains 12-five remembered rings\\
C) Reaction Ni$_{\text{(impure)}}$ + 4CO $\xrightarrow{450-470\text{K}}$ Ni(CO)$_4$ $\xrightarrow{330-350\text{K}}$ Ni$_{\text{(pure)}}$ + 4CO is used to purify Ni metal\\
D) Hydration energy of Cu$^{2+}$ is independent of cfse of Cu$^{2+}$\\
E) As(C$_6$H$_5$)$_3$ can form d$\pi-$d$\pi$ bond with transition metal

\qn{9} Among the following, how many have only one nodal plane $\sigma_{1S}$, $\sigma_{2S}$, $\sigma^*_{1S}$, $\pi_{2px}$, $\pi^*_{2py}$

\qn{10} In how many reactions one of the product is obtained as black precipitate?\\
i) Bi(OH)$_3$ + $\left[\text{Sn(OH)}_4\right]^{-2}_{\text{(aq.)}}\rightarrow$\\
ii) Bi$^{+3}_{\text{(aq.)}}$ + I$^-_{\text{(aq)}}$ (not excess) $\rightarrow$\\
iii) Bi$^{+3}_{\text{(aq.)}}$ + H$_2$S$_{\text{(g)}}$ $\xrightarrow{\text{H}^+}$\\
iv) $\left[\text{BiI}_4\right]^{-1}_{\text{(aq.)}}$ + H$_2$O $\xrightarrow{\text{dilution}}$\\
v) Bi$^{+3}_{\text{(aq.)}}$ + C$_9$H$_7$ON (5\%) + H$^+_{\text{(aq)}}$ + I$^-$(6M) $\rightarrow$\\
vi) Bi$^{+3}_{\text{(aq.)}}$ + C$_6$H$_3$(OH)$_3$(10\%) $\rightarrow$\\
vii) BiO.OH + H$_2$O$_{2\text{(aq.)}}$ $\rightarrow$

\partsec{3}{MATHEMATICS}
\sect{Section-I(i)}

\qn{1} Let $f:[a,b]\to[a,b]$ (where $a<b$ and are real numbers) be a differentiable non-constant function for which $f(a)=b$ and $f(b)=a$. Then
\begin{qoptions}
\opt{A}{$\exists$ at least one $c\in(a,b)$ such that $f(c)=c$}
\opt{B}{$\exists\ c_1,c_2\in(a,b),\ c_1\neq c_2$ such that $f'(c_1).f'(c_2)=1$}
\opt{C}{$\exists\ c\in(a,b)$ such that $|f'(c)|>1$}
\opt{D}{$\exists\ c\in(a,b)$ such that $|f'(c)|<1$}
\end{qoptions}

\qn{2} Let A be an $n\times n$ Nilpotent Matrix with real entries, then
\begin{qoptions}
\opt{A}{$A^k=0$ for some $k\in\mathbb N$}
\opt{B}{$I+A$ is invertible}
\opt{C}{$\det(I+A)$ is constant and independent of elements of A}
\opt{D}{$\det(I+A)>0$}
\end{qoptions}

\qn{3} Let $f:(0,\infty)\to\mathbb R$, has an anti-derivative g such that $f(x).g\left(\dfrac1x\right)=\dfrac1x\ \forall x>0$ and $f(1)=2$ Then
\begin{qoptions}
\opt{A}{$\left\lfloor f\left(\dfrac12\right)\right\rfloor=0$}
\opt{B}{$\left\lfloor\sqrt{f(2)}\right\rfloor=4$}
\opt{C}{$\lfloor f(1)\rfloor=2$}
\opt{D}{$\lfloor f(1)\rfloor=1$}
\end{qoptions}

\qn{4} If $x,y\in(-2,2)$ and $xy=-1$. Also, $u=\dfrac4{4-x^2}+\dfrac9{9-y^2}$. Then,
\begin{qoptions}
\opt{A}{minimum value of u is $\dfrac{67}{25}$}
\opt{B}{minimum value of u is $\dfrac{12}5$}
\opt{C}{u will be minimum when $x=\sqrt{\dfrac23}$}
\opt{D}{u will be minimum when $x=\sqrt{\dfrac32}$}
\end{qoptions}

\qn{5} Contents of the two urns is as given in this table. A fair die is tossed. If the face 1, 2, 4 or 5 comes, a marble is drawn from the urn A otherwise a marble is chosen from the urn B.\par\nobreak
\begin{center}\small
\begin{tabular}{|c|c|c|c|}
\hline
Urn & Red marbles & White marbles & Blue marbles \\
\hline
A & 5 & 3 & 8 \\
\hline
B & 3 & 5 & 0 \\
\hline
\end{tabular}
\end{center}
Let $E_1$: Denote the event that a red marble is chosen.\\
$E_2$: Denotes the events that a white marble is chosen\\
$E_3$: Denote the event that a blue is chosen.\\
Then
\begin{qoptions}
\opt{A}{Events $E_1$, $E_2$ and $E_3$ are equiprobable.}
\opt{B}{$P(E_1)$, $P(E_2)$, $P(E_3)$ are in A.P.}
\opt{C}{If the marble drawn is red, the probability that it came from the urn A is $\dfrac12$.}
\opt{D}{If the marble drawn is white, the probability that the face 5 appeared on the die is $\dfrac3{32}$.}
\end{qoptions}

\qn{6} Let $A\in M_2(\mathbb C)$ and $A^2=-I_2$. Then we can have
\begin{qoptions}
\opt{A}{$A=iI_2$}
\opt{B}{$A=-iI_2$}
\opt{C}{A is a periodic matrix}
\opt{D}{Number of solution for A is 2}
\end{qoptions}

\sect{Section-I(ii)}

\paraq{7}{8}
Let O be an interior point of $\Delta ABC$ such that $\overrightarrow{OA}+2\overrightarrow{OB}+3\overrightarrow{OC}=\vec0$. Point D \& E are the midpoints of the sides AC and BC respectively.

\qn{7} $\dfrac{\text{ar}(\Delta\,DOE)}{\text{ar}(\Delta\,ABC)}=$
\begin{qoptions}
\opt{A}{$\dfrac13$}
\opt{B}{$\dfrac16$}
\opt{C}{$\dfrac19$}
\opt{D}{None of these}
\end{qoptions}

\qn{8} $\dfrac{\text{ar}(\Delta\,ABC)}{\text{ar}(\Delta\,AOC)}=$
\begin{qoptions}
\opt{A}{2}
\opt{B}{3}
\opt{C}{4}
\opt{D}{8}
\end{qoptions}

\paraq{9}{10}
Let $a_0,a_1,a_2,\dots\dots\dots,a_n,\dots\dots\dots\dots$ be a sequence of numbers satisfying $(3-a_{n+1})(6+a_n)=18$ and $a_0=3$.

\qn{9} $a_5=$
\begin{qoptions}
\opt{A}{$\dfrac13$}
\opt{B}{$\dfrac17$}
\opt{C}{$\dfrac1{21}$}
\opt{D}{$\dfrac1{63}$}
\end{qoptions}

\qn{10} $\displaystyle\sum_{i=0}^{8}\dfrac1{a_i}=$
\begin{qoptions}
\opt{A}{$\dfrac{1013}3$}
\opt{B}{$\dfrac{1010}3$}
\opt{C}{100}
\opt{D}{None of these}
\end{qoptions}

\paraq{11}{12}
Consider three points $A\left(0,\dfrac43\right)$, $B(-1,0)$ and $C(1,0)$. The distance from a variable point P to the line BC is the geometric mean of the distances from this point to the lines AB and AC.

\qn{11} The locus of point P contains
\begin{qoptions}
\opt{A}{a circle \& a parabola}
\opt{B}{a parabola \& an ellipse}
\opt{C}{a circle \& a hyperbola}
\opt{D}{a hyperbola \& an ellipse}
\end{qoptions}

\qn{12} If a line L of zero slope passes through the Incenter of $\Delta ABC$, then number of distinct points of intersection with locus of P is :\underline{\hspace{1cm}}
\begin{qoptions}
\opt{A}{1}
\opt{B}{2}
\opt{C}{3}
\opt{D}{4}
\end{qoptions}

\sect{Section-III}

\qn{1} $f:\mathbb R\setminus\{0,1\}\to\mathbb R$ such that $f(x)+2f\left(\dfrac1x\right)+3f\left(\dfrac x{x-1}\right)=x$, then $16f(4)=\dots\dots\dots\dots$

\qn{2} $\displaystyle\lim_{n\to\infty}n\left(e-\left(1+\dfrac1n\right)^n\right)=a$ and $\displaystyle\lim_{n\to\infty}n\left(\left(1+\dfrac1n\right)^{n+1}-e\right)=b$ then $\lfloor 2a+2b\rfloor=\dots\dots\dots\dots$

\qn{3} Let $F:\mathbb R\to\mathbb R$, anti derivative of the function $f:\left[\dfrac\pi3,\dfrac{2\pi}3\right]\to\mathbb R$, $f(x)=\dfrac x{\sin x}$; and $F\left(\dfrac{2\pi}3\right)-F\left(\dfrac\pi3\right)=\dfrac\pi2\ln c$. Then $c=\dots\dots\dots\dots\dots\dots$

\qn{4} If $x,y\in\mathbb R$ and $\log_4(x+2y)+\log_4(x+2y)=1$. Then the minimum value of $(|x|-|y|)^2$ is $\dots\dots\dots\dots$

\qn{5} As shown in the diagram, points $P_1,P_2,\dots\dots\dots P_{10}$ are either the vertices or the midpoints of the edge of a tetrahedron respectively. Then, there are N number of groups of four points $(P_i,P_j,P_k,P_l)$; $(1\le i<j<k\le 10)$ such that in a group all four points are on the same plane. The sum of digits of N is $\dots\dots\dots$
\qfig{p2_math_s3_q05_1}

\qn{6} If $\displaystyle\lim_{x\to\infty}\dfrac{1^n+2^n+3^n+\dots\dots+n^n}{n^n}=a$ then $\left\lfloor\dfrac{10}{a}\right\rfloor$ is $\dots\dots\dots\dots\dots$

\qn{7} The number of real solutions of equation $2017^x-2016^x=x$ is\underline{\hspace{1.5cm}}

\qn{8} Consider $f(x)=ax^2+bx+c$; $a,b,c\in R$ and $a\neq0$. If satisfies the following conditions;\\
(i) $\forall\ x\in R,\ f(x-4)=f(2-x)$ and $f(x)\ge x$\\
(ii) $\forall\ x\in(0,2),\ f(x)\le\left(\dfrac{x+1}2\right)^2$\\
(iii) the minimum value of $f(x)$ on R is zero.\\
Now $\forall\ x\in[1,m]$, $f(x+t)\le x$ for some fixed $t\in R$ where $m>1$, the maximum `m' is $\dots\dots\dots\dots$

\qn{9} If $\left(1+\dfrac1n\right)^{n+x_n}=e$, $n\geqslant1$ $\displaystyle\lim_{n\to\infty}x_n=a$, then unit digit in $\left\lfloor\dfrac2a\right\rfloor$ is $\dots\dots\dots$

\qn{10} A and B alternatively throw a fair coin. The player who throws head two times first will win. A has the first throw. If chance of winning of A is $\dfrac ab$; $a,b\in\mathbb N$, $\gcd(a,b)=1$, then unit digit of $a+b$ is

\end{multicols}
\end{document}
